\chapter{Lista de materiales}

\section{Componentes esenciales}

Los siguientes componentes son indispensables para el sistema en su configuración mínima (un solo teléfono). Precios referenciales en pesos chilenos a mediados de 2026.

\begin{longtable}{p{0.4\textwidth} p{0.32\textwidth} R{0.18\textwidth}}
\toprule
\textbf{Componente} & \textbf{Modelo sugerido} & \textbf{Precio CLP} \\
\midrule
\endhead
Single Board Computer & Raspberry Pi 4 Modelo B (2\,GB) & \$85.000 \\
Fuente de poder & USB-C oficial Raspberry Pi 15\,W (5\,V/3\,A) & \$15.000 \\
Almacenamiento & MicroSD SanDisk 32\,GB Clase 10 / A1 & \$10.000 \\
Disipación térmica & Caja con ventilador para Pi 4 & \$12.000 \\
Cable de video & Cable micro-HDMI a HDMI, 1.5\,m & \$5.000 \\
Adaptador telefónico & Grandstream HT801 (1 puerto FXS) & \$40.000 \\
Cable telefónico & Cable RJ11, 2\,m & \$2.000 \\
Cable de red & Patch cord Ethernet Cat5e, 1\,m & \$3.000 \\
Teléfono analógico & Cualquier teléfono fijo con DTMF (años 90+) & --- \\
Televisor/monitor & Equipo con entrada HDMI & --- \\
\midrule
\textbf{Total esenciales} & & \textbf{\textasciitilde\$172.000} \\
\bottomrule
\caption{Bill of materials esencial.}\\
\end{longtable}

\subsection{Detalles por componente}

\subsubsection{Raspberry Pi 4 (2\,GB)}

Disponible en Chile a través de distribuidores oficiales como \texttt{raspberrypi.cl} y \texttt{mcielectronics.cl}. Se requiere el modelo Pi 4B con al menos 2\,GB de RAM. Modelos anteriores (3B+ o inferiores) no tienen suficiente rendimiento de emulación, y modelos con menos RAM dejan sin margen al sistema.

\begin{nota}
Si la Pi 4 (2\,GB) está sin stock, la Pi 5 (2\,GB) sirve también y rinde mejor. La diferencia de precio suele ser de 10\% adicional. Cualquiera de las dos cumple los requisitos.
\end{nota}

\subsubsection{Fuente USB-C oficial}

\textbf{Crítica.} La Pi 4 demanda 5\,V a 3\,A reales bajo carga. Cargadores de teléfono móvil, incluso los de marca, suelen entregar voltaje inestable, lo cual provoca \emph{undervoltage warnings}, reinicios espontáneos y corrupción de la microSD. La fuente oficial es la única forma confiable de evitarlo.

\subsubsection{MicroSD}

32\,GB es más que suficiente. Las especificaciones importantes son:

\begin{itemize}
  \item \textbf{Clase 10}: garantiza 10\,MB/s de escritura sostenida.
  \item \textbf{Endurance A1 o A2}: optimizado para aplicaciones (no solo grandes archivos como video).
\end{itemize}

Marcas recomendadas: SanDisk Endurance, Samsung Evo Plus.

\begin{tip}
Las microSD ``baratas sin marca'' fallan en cuestión de meses cuando un sistema operativo escribe constantemente. Para un \emph{appliance} que vive encendido, vale la pena gastar un poco más.
\end{tip}

\subsubsection{Grandstream HT801}

\emph{Analog Telephone Adapter} (ATA) con un puerto FXS \emph{(Foreign eXchange Subscriber)}, es decir, un puerto que provee voltaje de línea (\textasciitilde48\,V DC) y voltaje de timbre (\textasciitilde90\,V AC) hacia un teléfono analógico, y por el otro lado se conecta por Ethernet a la red local y se registra como cliente SIP.

Características relevantes:

\begin{itemize}
  \item Codecs G.711 (PCMU/PCMA), G.722, G.729, iLBC, OPUS
  \item DTMF in-band, RFC 2833 y SIP INFO
  \item Configurable vía web
  \item Función \emph{Off-hook Auto-Dial}: marca automáticamente un número al descolgar
\end{itemize}

En Chile se consigue por MercadoLibre, distribuidores VoIP como Inttelec o Centralinet, o importación directa por Aliexpress (con riesgo de garantía).

\subsubsection{Teléfono analógico}

Cualquier teléfono fijo con teclado DTMF de los últimos cuarenta años funciona. \textbf{Los teléfonos de disco rotatorio no sirven} porque generan pulsos, no tonos. Modelos clásicos que cumplen el rol y aportan estética noventera:

\begin{itemize}
  \item Panasonic KX-T series
  \item General Electric con cable rizado
  \item VTech inalámbrico (los primeros, con base grande)
  \item Cualquier teléfono ``de hotel'' Cetis o similar
\end{itemize}

\section{Componentes opcionales}

\begin{longtable}{p{0.4\textwidth} p{0.4\textwidth} R{0.15\textwidth}}
\toprule
\textbf{Componente} & \textbf{Razón} & \textbf{Precio CLP} \\
\midrule
\endhead
TV CRT pequeño usado & Estética auténtica años 90 & \$20.000--40.000 \\
Switch Ethernet 5 puertos & Si el router está lejos del setup & \$10.000 \\
Capturadora HDMI USB & Para grabar/transmitir partidas & \$25.000 \\
Parlantes externos 3.5\,mm & Si el TV tiene audio pobre & \$15.000 \\
Etiquetas decorativas & Diseño tipo ``Jugador 1'' en el teléfono & \$3.000 \\
\bottomrule
\caption{Componentes opcionales.}\\
\end{longtable}

\section{Plan de crecimiento: multijugador}

El sistema base soporta un solo teléfono. Para escalar al modo televisión con múltiples jugadores simultáneos:

\begin{itemize}
  \item Reemplazar el HT801 por un \textbf{HT802} (2 puertos FXS, $\sim$\$55.000) o \textbf{HT814} (4 puertos FXS, $\sim$\$110.000).
  \item Adquirir teléfonos adicionales.
  \item Reconfigurar Asterisk con \texttt{ConfBridge} y lógica de cola de turnos.
\end{itemize}

Este crecimiento se documenta en detalle en el capítulo 9.

\section{Comparativa con alternativas}

\begin{table}[h]
\centering
\begin{tabular}{lccc}
\toprule
& \textbf{Pi 4 (2\,GB)} & \textbf{Pi 5 (2\,GB)} & \textbf{Mini PC x86 usado} \\
\midrule
Precio aprox. (CLP) & \$85.000 & \$95.000 & \$100.000 \\
Emulación de Hugo & Buena & Muy buena & Excelente (nativo) \\
Consumo eléctrico & 5--7\,W & 8--12\,W & 15--25\,W \\
Asterisk & Perfecto & Perfecto & Perfecto \\
Tamaño físico & Tarjeta de crédito & Tarjeta de crédito & Pequeño \\
HDMI a TV & Sí (micro) & Sí (micro) & Sí (estándar) \\
Estética appliance & Excelente & Excelente & Buena \\
\bottomrule
\end{tabular}
\caption{Comparativa de plataformas para Hugo Phone.}
\end{table}

Para este proyecto se selecciona la \textbf{Pi 4 (2\,GB)} por la combinación de precio, tamaño y eficiencia energética.
